% Analisis de los resultados de rendimiento (se incluye desde main.tex).
% Si se vuelve a ejecutar el benchmark, revise que las cifras citadas aqui
% sigan correspondiendo con las tablas.

\paragraph{Sobrecosto de la paralelización.} Con un hilo, las versiones OpenMP y Pthreads tardan
prácticamente lo mismo que la secuencial: las diferencias son menores al 5\,\% y están dentro del ruido de
medición. Esto confirma que crear una única región paralela y llamar a los kernels por filas no añade un
costo apreciable.

\paragraph{Escalabilidad.} Para pocos hilos la aceleración es buena: con 2 hilos se obtiene un speedup de
1.6--1.9 y con 4 hilos de 2.6--3.4. El mejor resultado se alcanza alrededor de \textbf{8 hilos}, con un
speedup de aproximadamente \textbf{4.7} en el mundo de $1000\times1000$, 4.2 en el de $2000\times2000$ y 3.5
en el de $500\times500$. A partir de 8--12 hilos el tiempo deja de mejorar, y con 24--32 hilos incluso
empeora. Varias causas contribuyen a este comportamiento:
\begin{itemize}[nosep]
  \item \textbf{Kernel limitado por memoria.} Cada paso hace muy poco cómputo por byte transferido:
    compara tipos de celdas vecinas y copia el estado de un buffer a otro. Un microbenchmark de copia de
    arreglos en esta misma máquina deja de escalar a partir de 8 hilos, igual que la simulación.
  \item \textbf{Hilos lógicos (SMT).} Con más de 16 hilos, dos hilos comparten un mismo núcleo físico y sus
    recursos de caché y de carga/almacenamiento, lo que no aporta beneficio en un código limitado por memoria.
  \item \textbf{Barreras.} Cada generación tiene 4 barreras. En el mundo de $500\times500$ una generación
    dura apenas unos milisegundos, así que el costo de sincronizar 16--32 hilos pesa proporcionalmente más.
    Por eso es el caso que peor escala.
  \item \textbf{Plataforma.} Las mediciones se hicieron en un portátil con Windows y el plan de energía
    ``Equilibrado''. La frecuencia del procesador baja cuando se activan muchos núcleos, y el planificador de
    Windows migra hilos entre los dos complejos de núcleos (CCD) del Ryzen. Incluso un bucle puramente de
    cómputo, sin acceso a memoria, alcanzó solo $\approx6\times$ con 16 hilos en este equipo. Esto también
    explica la variabilidad observada; por ejemplo, algunas configuraciones con 12 o 24 hilos son
    notablemente más lentas que sus vecinas.
\end{itemize}

\paragraph{Planificación.} Las tres planificaciones de OpenMP rinden de forma muy parecida, sin una
ganadora consistente. Es lo esperado: en el esquema \emph{decidir--recoger} cada fila tiene un costo casi
constante, porque se recorren todas las celdas sin importar cuántos animales haya, de modo que el reparto
\texttt{static} ya está balanceado. Las planificaciones adaptativas (\texttt{dynamic}, \texttt{guided}) solo
compensarían si la densidad de animales hiciera muy desigual el trabajo por fila. Por otro lado, tienen la
desventaja de que un hilo puede recibir filas distintas en cada paso, lo que empeora la localidad de caché.
Pthreads, con bloques estáticos de filas, se comporta como OpenMP \texttt{static} hasta 8 hilos. Con más
hilos sufre algo más, porque su barrera (\texttt{pthread\_barrier} de winpthreads) es menos eficiente que la
de libgomp.

\paragraph{Tamaño del problema.} El mundo de $1000\times1000$ es el que mejor escala: su estado de unos
11\,MB cabe en la caché L3. El de $2000\times2000$ (unos 44\,MB) excede la caché y depende más del ancho de
banda de la memoria principal, y en el de $500\times500$ domina el costo de las barreras.
